\documentclass[9pt,landscape,a4paper]{article}
\usepackage[utf8]{inputenc}
\usepackage[T1]{fontenc}
\usepackage{geometry}
\usepackage{multicol}
\usepackage{amsmath,amssymb}
\usepackage{enumitem}
\usepackage{graphicx}
\usepackage{xcolor}
\usepackage{listings}
\usepackage{booktabs}
\usepackage{titlesec}
\usepackage{fancyhdr}
\usepackage{hyperref}
\usepackage{CJKutf8}

% Page layout
\geometry{top=0.7cm, bottom=0.7cm, left=0.7cm, right=0.7cm}
\pagestyle{empty}
\setlength{\parindent}{0pt}
\setlength{\parskip}{1.5pt}
\setlength{\columnsep}{7pt}

% Section formatting
\titleformat{\section}{\normalsize\bfseries\color{blue!70!black}}{}{0em}{}[\titlerule]
\titleformat{\subsection}{\small\bfseries\color{blue!50!black}}{}{0em}{}
\titlespacing{\section}{0pt}{4pt}{2pt}
\titlespacing{\subsection}{0pt}{3pt}{1pt}

% Code listing style
\lstset{
    language=Python,
    basicstyle=\ttfamily\fontsize{6.5}{7.5}\selectfont,
    keywordstyle=\color{blue!70!black}\bfseries,
    stringstyle=\color{red!60!black},
    commentstyle=\color{green!50!black}\itshape,
    numberstyle=\tiny\color{gray},
    backgroundcolor=\color{gray!8},
    frame=single,
    framerule=0.3pt,
    rulecolor=\color{gray!50},
    breaklines=true,
    breakatwhitespace=false,
    tabsize=2,
    showstringspaces=false,
    aboveskip=2pt,
    belowskip=2pt,
    xleftmargin=2pt,
    xrightmargin=2pt,
}

\newcommand{\keybox}[1]{%
    \begin{center}
    \fcolorbox{blue!50!black}{blue!5}{%
        \parbox{0.95\linewidth}{\small #1}%
    }
    \end{center}
}

\newcommand{\zhint}[1]{{\scriptsize\color{gray!70!black}#1}}

\begin{document}
\begin{CJK}{UTF8}{gbsn}

\begin{center}
{\large\bfseries\color{blue!70!black} Chapter 6 -- Data Modeling I}\\
{\footnotesize IERG4320/IEMS5726 Data Science in Practice | Cheat Sheet}
\end{center}

\begin{multicols}{3}

% ==================== SECTION 1 ====================
\section{Machine Learning Basics}
\zhint{ML = 让计算机从数据中学习规律，而不是人工写规则}

\subsection{ML Types}
\begin{itemize}[leftmargin=8pt,itemsep=0pt,topsep=0pt]
\item \textbf{Supervised}: labeled data $\to$ Regression / Classification \zhint{有标签}
\item \textbf{Unsupervised}: no labels $\to$ Clustering / Dim. Reduction \zhint{无标签，找规律}
\item \textbf{Semi-supervised}: small labeled + large unlabeled \zhint{少量标签+大量无标签}
\item \textbf{Reinforcement}: agent learns via rewards in environment \zhint{强化学习，通过奖励学}
\end{itemize}

\subsection{Regression vs Classification}
\textbf{Regression}: predict continuous value (price, count)\\
\zhint{回归：预测连续值，如房价}\\
\textbf{Classification}: predict discrete class (spam/not, cat/dog)\\
\zhint{分类：预测离散类别，如垃圾邮件检测}

% ==================== SECTION 2 ====================
\section{Classification Algorithms}

\subsection{1. Decision Tree}
\zhint{决策树：用属性做判断，像if-else树状结构}\\
\textbf{Key}: uses Gini Index or Entropy to split nodes.\\
Lower Gini = better split (less misclassification).

\begin{lstlisting}
from sklearn.tree import DecisionTreeClassifier
from sklearn.model_selection import train_test_split
from sklearn import metrics
# Load data
X = pima[feature_cols]  # Features
y = pima.Outcome        # Target
X_train, X_test, y_train, y_test = \
  train_test_split(X, y, test_size=0.3,
                   random_state=1)
# Build and evaluate
model = DecisionTreeClassifier()
model = model.fit(X_train, y_train)
y_pred = model.predict(X_test)
print("Accuracy:",
      metrics.accuracy_score(y_test, y_pred))
\end{lstlisting}

\subsection{2. Random Forest}
\zhint{随机森林：多棵决策树投票，减少过拟合}\\
Average multiple decision trees trained on different data subsets. Final output = majority vote.

\begin{lstlisting}
from sklearn.ensemble import RandomForestClassifier
rf = RandomForestClassifier(
  n_estimators=50, max_features="sqrt",
  random_state=44)
rf.fit(X_train, y_train)
preds = rf.predict(X_test)
print(metrics.accuracy_score(y_test, preds))
\end{lstlisting}
\zhint{n\_estimators=树的数量，max\_features=每次分裂用的特征数}

\subsection{3. SVM (Support Vector Machine)}
\zhint{SVM：找最大间隔的超平面来分类}\\
Find the hyperplane with maximum margin. Support vectors = data points closest to boundary.\\
\textbf{Kernels}: linear, poly, rbf (for non-linear data).

\begin{lstlisting}
from sklearn import svm, datasets, metrics
cancer = datasets.load_breast_cancer()
X_train, X_test, y_train, y_test = \
  train_test_split(cancer.data,cancer.target,
                   test_size=0.3,random_state=109)
clf = svm.SVC(kernel='linear')
clf.fit(X_train, y_train)
y_pred = clf.predict(X_test)
print("Accuracy:",
      metrics.accuracy_score(y_test, y_pred))
# Accuracy: 0.9649
\end{lstlisting}

\subsection{4. KNN (K-Nearest Neighbor)}
\zhint{KNN：找最近的K个邻居，多数投票决定类别}\\
Classify based on majority class of K nearest neighbors. Distance: Euclidean.

\begin{lstlisting}
from sklearn.neighbors import KNeighborsClassifier
from sklearn.datasets import load_iris
irisData = load_iris()
X_train,X_test,y_train,y_test = \
  train_test_split(irisData.data,irisData.target,
                   test_size=0.2,random_state=42)
knn = KNeighborsClassifier(n_neighbors=7)
knn.fit(X_train, y_train)
print("Accuracy:",
      metrics.accuracy_score(y_test,
                             knn.predict(X_test)))
# Accuracy: 0.9667
\end{lstlisting}

\subsection{5. ANN (Artificial Neural Network)}
\zhint{人工神经网络：模拟神经元连接，可建模非线性关系}

\begin{lstlisting}
import torch
import torch.nn as nn
n_input, n_hidden, n_out = 10, 15, 1
batch_size, learning_rate = 100, 0.01
data_x = torch.randn(batch_size, n_input)
data_y = (torch.rand(size=(batch_size,1))<0.5)\
         .float()
# Define model
model = nn.Sequential(
  nn.Linear(n_input, n_hidden), nn.ReLU(),
  nn.Linear(n_hidden, n_out), nn.Sigmoid())
loss_function = nn.MSELoss()
optimizer = torch.optim.SGD(
  model.parameters(), lr=learning_rate)
# Training loop
losses = []
for epoch in range(50000):
  pred_y = model(data_x)
  loss = loss_function(pred_y, data_y)
  losses.append(loss.item())
  model.zero_grad()
  loss.backward()   # Backpropagation
  optimizer.step()  # Gradient descent
\end{lstlisting}
\zhint{ReLU隐藏层 + Sigmoid输出层(二分类); MSELoss + SGD优化}

% ==================== SECTION 3 ====================
\section{Classification Metrics}
\zhint{评估指标：衡量模型预测好不好}

\subsection{Accuracy}
$\text{Accuracy} = \frac{\text{correct predictions}}{\text{total predictions}}$\\
\zhint{准确率 = 预测对的/总数。缺点：类别不平衡时不可靠}

\begin{lstlisting}
from sklearn import metrics
metrics.accuracy_score(y_test, y_pred)
\end{lstlisting}

\subsection{Confusion Matrix}
\zhint{混淆矩阵：TP/TN/FP/FN四个值}
\begin{center}
\scriptsize
\begin{tabular}{c|c|c}
 & Pred Pos & Pred Neg \\ \hline
Actual Pos & TP \zhint{真阳} & FN \zhint{漏报} \\ \hline
Actual Neg & FP \zhint{误报} & TN \zhint{真阴} \\
\end{tabular}
\end{center}

\begin{lstlisting}
from sklearn.metrics import confusion_matrix
m = confusion_matrix(y_test, y_pred)
\end{lstlisting}

\subsection{Precision \& Recall}
$\text{Precision} = \frac{TP}{TP + FP}$ \zhint{预测为正中，真正为正的比例}\\[2pt]
$\text{Recall} = \frac{TP}{TP + FN}$ \zhint{实际为正中，被正确找出的比例}\\[2pt]
\textbf{Trade-off}: high threshold $\to$ high precision, low recall; low threshold $\to$ opposite.\\
\zhint{阈值高=保守(精确但漏检多)，阈值低=激进(召回高但误报多)}

\subsection{F1 Score}
$F_1 = \frac{2 \cdot \text{Precision} \cdot \text{Recall}}{\text{Precision} + \text{Recall}}$
\zhint{F1是精确率和召回率的调和平均}\\
\textbf{Flaw}: F1 is class-variant (flipping pos/neg changes F1).

\subsection{MCC (Matthews Correlation Coefficient)}
\zhint{MCC：对不平衡数据更公平，范围[-1, 1]}\\
$\text{MCC} = \frac{TP \cdot TN - FP \cdot FN}{\sqrt{(TP+FP)(TP+FN)(TN+FP)(TN+FN)}}$\\
\textbf{Advantage}: class-invariant, range $[-1, 1]$.

\begin{lstlisting}
metrics.matthews_corrcoef(y_test, y_pred)
metrics.f1_score(y_test, y_pred)
\end{lstlisting}

\subsection{AUC-ROC}
\zhint{ROC曲线下面积：衡量模型区分正负样本的能力}\\
ROC = plot TPR vs FPR at different thresholds.\\
$\text{TPR} = \frac{TP}{TP+FN}$, \quad $\text{FPR} = \frac{FP}{FP+TN}$\\
AUC = area under ROC curve. AUC=0.5 = random guess, AUC=1 = perfect.

% ==================== SECTION 4 ====================
\section{Model Comparison (10-fold CV)}
\zhint{10折交叉验证：公平比较不同算法}

\begin{lstlisting}
from sklearn.model_selection import KFold, \
  cross_val_score
from sklearn.linear_model import LogisticRegression
from sklearn.tree import DecisionTreeClassifier
from sklearn.neighbors import KNeighborsClassifier
from sklearn.ensemble import RandomForestClassifier
from sklearn.naive_bayes import GaussianNB
from sklearn.svm import SVC

models = []
models.append(('LR',
  LogisticRegression(solver="liblinear")))
models.append(('KNN',
  KNeighborsClassifier(n_neighbors=5)))
models.append(('SVM',SVC(kernel='rbf')))
models.append(('NB', GaussianNB()))
models.append(('DT',
  DecisionTreeClassifier(criterion="entropy")))
models.append(('RF',
  RandomForestClassifier(n_estimators=100,
                     criterion="entropy")))
for name, model in models:
  kfold = KFold(n_splits=10)
  result = cross_val_score(model, X, y,
    cv=kfold, scoring='accuracy')
  print("%s: Mean=%.2f%% SD=%.2f%%" %
    (name,result.mean()*100,result.std()*100))
\end{lstlisting}

\subsection{Hyperparameter Tuning}
\zhint{超参数调优：选最佳模型设置}\\
\textbf{Parameters} (learned): weights, biases, centroids\\
\textbf{Hyperparameters} (set by user): learning rate, \# layers, dropout, epochs, batch size, kernel\\[2pt]
\textbf{Tuning methods}:
\begin{enumerate}[leftmargin=10pt,itemsep=0pt,topsep=0pt]
\item \textbf{Grid Search}: try all combinations \zhint{网格搜索}
\item \textbf{Manual}: domain expert picks \zhint{专家经验}
\item \textbf{Random Search}: random combos \zhint{随机搜索}
\end{enumerate}
\zhint{不确定就用默认参数，不要只报告最好的一次}

% ==================== SECTION 5 ====================
\section{Regression}
\zhint{回归：预测连续值，分析变量间关系}

\subsection{Linear Regression}
$y = a + bX$ (simple) \zhint{一元线性回归}\\
$y = b_0 + b_1x_1 + b_2x_2 + \cdots + b_nx_n$ (multiple)\\
\zhint{多元线性回归：多个特征预测一个目标值}\\
Vector form: $y = \theta^T x$\\
Learning = find best $\theta$ via Ordinary Least Squares (OLS).

\begin{lstlisting}
from sklearn import linear_model
import statsmodels.api as sm
X = df[['Interest_Rate','Unemployment_Rate']]
Y = df['Stock_Index_Price']
regr = linear_model.LinearRegression()
regr.fit(X, Y)
print('Intercept:', regr.intercept_)
print('Coef:', regr.coef_)
# Predict
pred = regr.predict([[2.75, 5.3]])
\end{lstlisting}

\subsection{R$^2$ and Adjusted R$^2$}
\zhint{R$^2$：模型拟合好坏，越接近1越好}\\
$R^2 = 1 - \frac{SS_{res}}{SS_{tot}}$ where $SS_{tot}=\sum(y_i-\bar{y})^2$, $SS_{res}=\sum(y_i-\hat{y}_i)^2$\\[2pt]
\textbf{Problem}: R$^2$ never decreases with more variables.\\
\zhint{Adjusted R$^2$修正了变量数量的影响}\\
$R^2_{adj} = 1 - (1-R^2)\frac{n-1}{n-k-1}$ \quad (k = \# of features)

\begin{lstlisting}
r2 = regr.score(X, Y)   # 0.8976
adj_r2 = 1-(1-r2)*(len(Y)-1)\
           /(len(Y)-X.shape[1]-1)  # 0.8879
\end{lstlisting}

\subsection{statsmodels OLS}
\zhint{statsmodels给出完整统计信息：p值、置信区间等}

\begin{lstlisting}
X_sm = sm.add_constant(X)  # add intercept
model = sm.OLS(Y, X_sm).fit()
print(model.summary())
# Shows: R2, Adj R2, F-statistic,
#        coef, std err, t, P>|t|, CI
\end{lstlisting}

\subsection{Regression Metrics}
\zhint{回归评估：MAE, RMSE, MSE}\\[2pt]
\textbf{MAE} = $\frac{1}{n}\sum|y_i - \hat{y}_i|$ \zhint{平均绝对误差}\\[2pt]
\textbf{MSE} = $\frac{1}{n}\sum(y_i - \hat{y}_i)^2$ \zhint{均方误差，对大偏差惩罚更重}\\[2pt]
\textbf{RMSE} = $\sqrt{\text{MSE}}$ \zhint{均方根误差}

\begin{lstlisting}
from sklearn.metrics import mean_squared_error,\
  r2_score, mean_absolute_error
print('MSE:', mean_squared_error(y_test, y_pred))
print('MAE:', mean_absolute_error(y_test, y_pred))
\end{lstlisting}

% ==================== SECTION 6 ====================
\section{Logistic Regression}
\zhint{逻辑回归：用sigmoid把线性输出转成0-1概率，用于分类}\\
Apply sigmoid to linear regression output:\\
$g(z) = \frac{1}{1 + e^{-z}}$ \quad where $z = \theta^T x$\\
Output bounded in $[0, 1]$; tends to 0 or 1.

\begin{lstlisting}
from sklearn.linear_model import LogisticRegression
clf = LogisticRegression(random_state=0)
clf.fit(X_train, Y_train)
Y_pred = clf.predict(X_test)
m = confusion_matrix(Y_test, Y_pred)
# ROC curve
from sklearn.metrics import roc_curve, \
  roc_auc_score
lr_probs = clf.predict_proba(X_test)[:, 1]
lr_auc = roc_auc_score(Y_test, lr_probs)
fpr, tpr, _ = roc_curve(Y_test, lr_probs)
\end{lstlisting}

% ==================== SECTION 7 ====================
\section{More Regression Models}

\subsection{Ridge Regression}
\zhint{岭回归：加L2正则化，防止过拟合}\\
Adds L2 penalty. $\alpha$ controls penalty weight.

\begin{lstlisting}
from sklearn.linear_model import Ridge
from sklearn.pipeline import make_pipeline
from sklearn.preprocessing import StandardScaler
pipeline = make_pipeline(
  StandardScaler(), Ridge(alpha=1.0))
pipeline.fit(X_train, y_train)
y_pred = pipeline.predict(X_test)
print('MSE:', mean_squared_error(y_test, y_pred))
print('R2:', r2_score(y_test, y_pred))
\end{lstlisting}

\subsection{LASSO Regression}
\zhint{LASSO：加L1正则化，可以做特征选择(系数变0)}

\begin{lstlisting}
from sklearn.linear_model import Lasso
model = Lasso(alpha=1.0)
model.fit(X, y)
yhat = model.predict([row])
\end{lstlisting}

\subsection{ANN Regression (PyTorch)}
\zhint{神经网络回归：输出层用ReLU(不是Sigmoid)，损失用MSE}

\begin{lstlisting}
n_input, n_hidden, n_out = 7, 5, 1
model = nn.Sequential(
  nn.Linear(n_input, n_hidden), nn.ReLU(),
  nn.Linear(n_hidden, n_out), nn.ReLU())
loss_function = nn.MSELoss()
optimizer = torch.optim.SGD(
  model.parameters(), lr=0.01)
for epoch in range(5000):
  pred_y = model(data_x)
  loss = loss_function(pred_y, data_y)
  model.zero_grad()
  loss.backward()
  optimizer.step()
\end{lstlisting}

\subsection{SVR (Support Vector Regressor)}
\zhint{支持向量回归：用核函数处理非线性回归}

\begin{lstlisting}
from sklearn.svm import SVR
sc_X, sc_y = StandardScaler(), StandardScaler()
X = sc_X.fit_transform(X)
y = sc_y.fit_transform(y)
regressor = SVR(kernel='rbf')
regressor.fit(X, y)
y_pred = regressor.predict(X_new)
\end{lstlisting}

\subsection{Parameters vs Hyperparameters}
\zhint{参数=模型学出来的; 超参数=人手动设的}
\begin{center}
\scriptsize
\begin{tabular}{l|l|l}
\textbf{Model} & \textbf{Params} & \textbf{Hyperparams} \\ \hline
MLR & coef, bias & None \\
Logistic & coef, bias & solver, penalty \\
Ridge & coef, bias & $\lambda$ (alpha) \\
LASSO & coef, bias, L1 & $\lambda$ (alpha) \\
\end{tabular}
\end{center}

\keybox{
\textbf{Regression Comparison}: Preprocess $\to$ Split 70/30 $\to$ Run MLR, Ridge, LASSO, SVR-rbf, ANN $\to$ Compare R$^2$, MAE, MSE.
\zhint{回归模型对比流程：和分类类似，用回归指标评估}
}

\end{multicols}
\end{CJK}
\end{document}
